% ============================================================
% 摘要 (Abstract) - 严格按照哈尔滨工程大学模板
% ============================================================

% ----- ABSTRACT -----
\clearpage
\section*{ABSTRACT}
\addcontentsline{toc}{section}{ABSTRACT}

With the continuous opening of Arctic shipping routes, safe navigation in polar ice-covered waters has become a critical challenge in naval architecture. This thesis presents a comprehensive framework that integrates Unreal Engine 5-based simulation, multi-modal data collection, 3D Gaussian Splatting (3DGS) reconstruction, and intelligent route planning for autonomous surface vessels operating in polar environments.

The research establishes a simulation platform using ASVSim (Autonomous Surface Vehicle Simulator) running on Unreal Engine 5, enabling realistic physics-based vessel dynamics and high-fidelity environmental rendering. A multi-modal data collection pipeline was developed to capture synchronized RGB images, depth maps, LiDAR point clouds, and instance segmentation masks at 640$\times$480 resolution. The system achieves an acquisition rate of approximately 6.5 seconds per frame through systematic optimization of rendering settings, including disabling Lumen global illumination and post-processing effects.

The methodology addresses key technical challenges including sensor configuration, data synchronization, and coordinate system alignment for 3DGS training. The collected dataset is formatted for COLMAP structure-from-motion and subsequent 3D Gaussian Splatting reconstruction. The proposed approach enables real-time ice condition perception and dynamic route re-planning, combining deep learning-based instance segmentation (SAM3) with metric depth estimation (Depth Anything 3).

Experimental validation demonstrates the effectiveness of the data collection system, with successful acquisition of 200-frame sequences containing synchronized multi-sensor data. The intelligent perception layer was deployed and evaluated: Depth Anything 3 achieved 11.36 FPS throughput processing 126 frames with metric depth accuracy (median error 2.67m after alignment), while SAM3 processed 94 frames detecting an average of 8.7 instances per frame with 0.947 mean confidence. Notably, analysis revealed DA3's pose prediction to be unreliable for texture-poor polar scenes, necessitating reliance on simulation ground truth poses. The results establish a foundation for future integration of perception, reconstruction, and planning modules in polar navigation scenarios.

\vspace{1em}
\noindent\textbf{Keywords:} Polar route planning; 3D Gaussian Splatting; Unreal Engine; Autonomous surface vessels; Multi-modal data collection; Computer vision; Naval architecture

% ----- 摘 要 -----
\clearpage
\section*{摘\quad 要}
\addcontentsline{toc}{section}{摘\quad 要}

随着北极航道的持续开放，极地冰区水域的安全航行已成为船舶与海洋工程领域的关键挑战。本论文提出了一个综合框架，集成基于 Unreal Engine 5 的仿真、多模态数据采集、3D Gaussian Splatting（3DGS）重建和智能路径规划，用于极地环境中作业的自主水面航行器。

本研究建立了基于 ASVSim（自主水面航行器仿真器）的仿真平台，该平台运行在 Unreal Engine 5 上，实现了基于物理的船舶动力学和高保真环境渲染。开发了多模态数据采集管道，以 640$\times$480 分辨率采集同步的 RGB 图像、深度图、LiDAR 点云和实例分割掩码。通过对渲染设置的系统性优化（包括禁用 Lumen 全局光照和后期处理效果），系统实现了约每帧 6.5 秒的采集速度。

研究方法解决了传感器配置、数据同步和 3DGS 训练的坐标系对齐等关键技术挑战。采集的数据集按照 COLMAP 运动恢复结构格式进行处理，用于后续的 3D Gaussian Splatting 重建。所提出的方法实现了实时冰情感知和动态路径重规划，结合了基于深度学习的实例分割（SAM3）和度量深度估计（Depth Anything 3）。

实验验证证明了数据采集系统的有效性，成功采集了包含同步多传感器数据的 200 帧序列。智能感知层已完成部署和评估：Depth Anything 3 处理 126 帧实现 11.36 FPS 的吞吐率，度量深度精度达 2.67 米（对齐后中位误差）；SAM3 处理 94 帧，平均每帧检测 8.7 个实例，平均置信度 0.947。值得注意的是，分析发现 DA3 的位姿预测在纹理贫乏的极地场景中不可靠，因此需依赖仿真真值位姿。这些结果为极地航行场景中感知、重建和规划模块的未来集成奠定了基础。

\vspace{1em}
\noindent\textbf{关键词：}极地路径规划；3D Gaussian Splatting；虚幻引擎；自主水面航行器；多模态数据采集；计算机视觉；船舶与海洋工程
